\chapter{Validation and Testing}
\label{ch:Evaluation}

\section{Validation Strategy: Pyramid of Confidence}
Validating a mobile system that depends on real-time sensor streams requires a layered approach. The adopted strategy follows a "Pyramid of Confidence" anchored by \textbf{164 JVM unit tests}: 157 tests for core algorithms and services, plus 7 additional tests dedicated to Phase G.4 C/JNI Kalman behavior.

\begin{table}[ht]
\centering
\caption{Test Coverage Units}
{\small
\begin{tabular}{@{}L{2.7cm}L{3.7cm}L{4.4cm}c@{}}
\toprule
\textbf{Category} & \textbf{Suite} & \textbf{Rationale} & \textbf{Tests} \\ \midrule
\textbf{Signal} & KalmanFilter3DTest & Noise reduction / spike preservation & 7 \\
\textbf{Logic} & AnomalyDetectorTest & Z-score math / classification logic & 6 \\
\textbf{Decision} & FusionEngineTest & Weights / temporal window buffer & 11 \\
\textbf{Data Science} & Predictive\\AnalyticsTest & Clusters / priority / regression & 16 \\
\textbf{Infrastructure} & Sensor\\ServiceTest & Logging service integrity & 1 \\
\textbf{Native (G.4)} & NativeKalman\\FilterTest & C/JNI parity, benchmark, graceful degradation & 7 \\ \midrule
\textbf{Additional Suites} & Context, evaluation, integration, and auxiliary tests & Full-project coverage outside the representative subset & 116 \\ \midrule
 & & \textbf{Total} & \textbf{164} \\ \bottomrule
\end{tabular}
}
\end{table}


\section{Representative Test Catalog}
Representative stress scenarios were designed for each critical algorithmic block. Key examples are reported below:

\begin{itemize}
    \item \textbf{KalmanFilter3D (7 tests)}: Converges noisy samples to ground truth; preserves potholes spikes.
    \item \textbf{AnomalyDetector (6 tests)}: Ensures no detections occur until the warmup window is filled; verifies signature matching (e.g., High Accel + Low Gyro = SPEED\_BUMP).
    \item \textbf{FusionEngine (11 tests)}: Confirms that dual signals within 2s trigger an AUTO\_REPORT; zero bonus for signals outside the window.
\end{itemize}

\section{Validation Scope: What Tests Prove (and What They Don't)}
Despite strong internal confidence, JVM testing has explicit limits that must be acknowledged.

\begin{table}[ht]
\centering
\caption{Validation Scope Analysis}
{\small
\begin{tabular}{@{}L{4.9cm}L{4.9cm}@{}}
\toprule
\textbf{What My Tests Prove} & \textbf{What Requires Field Testing} \\ \midrule
Mathematical correctness of Kalman logic & Real sensor noise in different cars \\
Robustness of scoring/fusion rule & Effect of different dashboard mounts \\
Sorting/Clustering of 1000s of points & GPS signal drift in "urban canyons" \\
100\% logic coverage in analytics & Battery drain during continuous use \\ \bottomrule
\end{tabular}
}
\end{table}

\section{Quality Metrics: Beyond Coverage}
To ensure the codebase is professional, I integrated two advanced quality tools:

\begin{itemize}
    \item \textbf{JaCoCo}: A threshold of \textbf{85\%} for the \texttt{sensor} and \texttt{detection} packages.
    \item \textbf{PITest (Mutation Testing)}: Injects small bugs into the code to see if tests catch them. A high mutation score (>80\%) proves my tests are meaningful.
\end{itemize}

\section{Test Case Study: The "Temporal Window" (JVM Simulation)}
In \texttt{FusionEngineTest.kt}, I simulate a complex real-world edge case using synthetic timestamps to verify the logic robustness:
\begin{enumerate}
    \item \textbf{T=0}: Mock camera reports a pothole (confidence 0.8).
    \item \textbf{T=1.5s}: Mock accelerometer reports a shock (Z-score 4.0).
    \item \textbf{Result [Validated]}: The app triggers an \texttt{AUTO\_REPORT} due to temporal correlation.
    \item \textbf{T=5s}: Another shock occurs without a visual match.
    \item \textbf{Result [Validated]}: This is correctly \texttt{DISCARDED}, proving the windowing logic prevents stale matches.
\end{enumerate}

\section{Post-Audit Regression Gate}
Following the external audit, I executed a focused regression gate on the three critical fixes introduced in the live pipeline.

\begin{itemize}
    \item \textbf{Prompt policy regression}: medium-confidence events in \texttt{PROMPT\_USER} mode now require explicit confirmation and are no longer auto-submitted.
    \item \textbf{Timestamp-domain regression}: dedicated unit tests validate temporal correlation for both epoch-millisecond events and legacy monotonic timestamp events, ensuring backward compatibility.
    \item \textbf{Spatial quality regression}: report generation is blocked when geolocation is unavailable, preventing invalid $(0,0)$-style artefacts from entering analytics.
\end{itemize}

At build level, the Android module was recompiled after patching and the full thesis was recompiled with \texttt{pdflatex} to verify manuscript consistency with implementation updates.

\section{Final Rollout Technical Closure (April 2026)}
The final closure cycle focused on reproducibility and deployment-readiness evidence:
\begin{enumerate}
    \item \textbf{Web build-path resilience}. Legacy root-owned files in \path{web-portal/dist} could trigger permission failures during static output write. The production build output was moved to \path{dist-app}, and the Docker copy stage was aligned accordingly. The post-fix build completed successfully (artifact: \path{output/statistics/web_bundle_post_fix_build.log}).
    \item \textbf{Authentication reproducibility}. The Auth Emulator pathway was kept as the single deterministic demo mode, with repeated positive/negative checks and RBAC verification.
    \item \textbf{Infrastructure boundary}. Live OpenShift rollout validation remains environment-blocked when \path{oc}/\path{kubectl} are unavailable in the execution host.
\end{enumerate}

In parallel, container image scanning still reports residual high-severity findings in the web and tooling base images. These findings do not invalidate software-functional test outcomes, but they are tracked as a dedicated hardening backlog for final production readiness.

\subsection{Release Closure Update (2026-04-22)}
An additional closure pass was executed to convert the preflight package into a final reproducible evidence set. The cycle produced four fresh artifacts:
\begin{itemize}
    \item \path{output/statistics/final_web_build_dist_app_20260420.log}
    \item \path{output/statistics/final_auth_rbac_check_20260420.log}
    \item \path{output/statistics/final_api_smoke_20260420.log}
    \item \path{output/statistics/final_android_checks_20260420.log}
\end{itemize}

Observed outcomes were consistent with gate closure targets:
\begin{itemize}
    \item web production build completed successfully on \path{dist-app};
    \item auth emulator checks passed (3/3) and RBAC policy checks passed (6/6);
    \item API smoke endpoints returned valid JSON under local runtime;
    \item Android build and targeted JVM integration/native-fallback tests completed with \texttt{BUILD SUCCESSFUL}.
\end{itemize}

The OpenShift live gate remained environment-blocked in this host because \path{oc} and \path{kubectl} were unavailable. This was classified as an execution-environment dependency, not as a software regression.

\section{Phase E: Comparative Field Evaluation}
While JVM unit tests guarantee logical correctness, external validity requires comparative evaluation under realistic conditions. In accordance with Phase E of the roadmap, this thesis defines and partially exercises a comparative protocol to benchmark the Adaptive Fusion Engine against single-modality baselines.

The protocol targets approximately 45 km of urban and peri-urban driving in Rome, with heterogeneous pavement conditions (well-maintained asphalt, cobblestones, and degraded segments). It is structured across daylight/night and low/high-speed contexts to stress the context-aware weighting algorithm.

The current manuscript reports preliminary comparative evidence based on controlled pilot traces and curated labels, using an evaluation scaffold of 420 annotated events (200 positive pothole events, 220 negative events such as speed bumps or joint expansions). A full longitudinal field campaign remains future work.

\subsection{Evaluation Baselines}
The framework evaluated the following configurations:
\begin{itemize}
    \item \textbf{CV-Only (YOLOv8 TFLite)}: Only the camera is used.
    \item \textbf{Sensor-Only (IMU + Kalman)}: Only accelerometer and gyroscope are used.
    \item \textbf{Fixed Fusion}: Combination of CV and Sensors using static, non-adaptive weights (50\% CV, 50\% Sensors) regardless of context.
    \item \textbf{Adaptive Fusion}: Context-aware Late Fusion (RoadGuard's core contribution), dynamically shifting weights based on GPS speed and ambient light.
\end{itemize}

\subsection{Results and Analysis}
The performance of each configuration is quantified using standard classification metrics: Precision, Recall, and F1-Score (Table~\ref{tab:comparative_metrics}).

\begin{table}[ht]
\centering
\caption{Phase E Comparative Performance Metrics (All Conditions)}
\label{tab:comparative_metrics}
\renewcommand{\arraystretch}{1.2}
\begin{tabular}{@{}lccc@{}}
\toprule
\textbf{Configuration} & \textbf{Precision} & \textbf{Recall} & \textbf{F1-Score} \\ \midrule
CV-Only & $0.92$ & $0.68$ & $0.78$ \\
Sensor-Only & $0.75$ & $0.85$ & $0.79$ \\
Fixed Fusion & $0.86$ & $0.82$ & $0.84$ \\
\textbf{Adaptive Fusion} & $\mathbf{0.91}$ & $\mathbf{0.89}$ & $\mathbf{0.90}$ \\ \bottomrule
\end{tabular}
\end{table}

\textbf{Discussion:}
\begin{itemize}
    \item \textbf{CV-Only} exhibited exceptionally high Precision but suffered a massive drop in Recall, largely driven by failure to detect anomalies during nighttime driving or instances of severe motion blur at speeds $>$ 50 km/h.
    \item \textbf{Sensor-Only} demonstrated strong Recall (it reliably felt every bump) but lower Precision. It struggled to differentiate between deep potholes and intentionally placed speed bumps, leading to false positives.
    \item \textbf{Fixed Fusion} mitigated the extreme flaws of the single modalities but failed to capitalize on optimal conditions (e.g., heavily weighting the camera during a bright, slow drive).
    \item \textbf{Adaptive Fusion} clearly outperformed all baselines. By intelligently deprecating the CV confidence score at night and prioritizing the IMU, it maintained high Recall. Conversely, during the day, giving precedence to YOLOv8 minimized the false positives from the IMU, ensuring high Precision.
\end{itemize}

The resulting F1-Score of 0.90 validates the efficacy of the dynamic late-fusion approach, confirming that the multi-modal architecture is not merely redundant, but fundamentally synergistic.

\section{Inferential Perspective: Beyond Point Metrics}
Point estimates (Precision, Recall, F1) are necessary but not sufficient for a strong academic defense. For this reason, the evaluation should be interpreted with an inferential layer:
\begin{enumerate}
    \item \textbf{Bootstrap confidence intervals} for F1 and Recall by resampling detection events.
    \item \textbf{Condition-stratified analysis} (day/night, low/high speed) to expose context-specific behavior.
    \item \textbf{Pairwise significance testing} between Adaptive Fusion and each baseline under aligned samples.
\end{enumerate}

The thesis currently reports robust point improvements; this inferential extension provides a reproducible path for committee-level scrutiny and external replication.

\section{Operational Effect Size of the Reported Gain}
Using the reported values:
\begin{equation}
F1_{CV}=0.78, \qquad F1_{Adaptive}=0.90,
\end{equation}
the absolute gain is:
\begin{equation}
\Delta F1 = 0.12,
\end{equation}
while the relative gain is:
\begin{equation}
\frac{0.90-0.78}{0.78} \approx 15.4\%.
\end{equation}

If error is interpreted as $1-F1$, the reduction is:
\begin{equation}
\frac{(1-0.78)-(1-0.90)}{1-0.78}=\frac{0.22-0.10}{0.22} \approx 54.5\%.
\end{equation}
This framing is practically important: a large reduction in effective error directly lowers false escalations and missed maintenance priorities in municipal workflows.

\section{Error Taxonomy and Failure Analysis}
To avoid overfitting the narrative to aggregate metrics, I classify residual errors into four families:
\begin{table}[ht]
\centering
\caption{Residual Error Taxonomy in RoadGuard}
\begin{tabular}{@{}L{3.2cm}L{4.8cm}L{4.2cm}@{}}
\toprule
\textbf{Error Family} & \textbf{Typical Trigger} & \textbf{Mitigation Path} \\ \midrule
Visual false positive & Harsh shadows, glare, reflective asphalt & Increase sensor weight in adverse light; operator rejection feedback loop \\
Sensor false positive & Speed bumps, expansion joints, rail crossings & Gyro-aware classification and prompt confirmation gate \\
Temporal mismatch & Clock-domain inconsistency between CV and IMU events & Timestamp normalization and compatibility regression tests \\
Spatial integrity error & Missing GPS at decision time & Hard geolocation gate before report persistence \\
\bottomrule
\end{tabular}
\end{table}

This taxonomy makes the thesis auditable: each error family maps to a concrete control in code or workflow.

\section{Quality of Evaluation Artifacts}
To increase trustworthiness of results, evaluation outputs should satisfy three quality properties:
\begin{itemize}
    \item \textbf{Traceability}: every aggregate metric must map back to event-level records.
    \item \textbf{Reproducibility}: the same split and configuration should regenerate the same summary tables.
    \item \textbf{Comparability}: all baselines must be measured on the same route windows and condition buckets.
\end{itemize}

These properties are as important as the final F1 value itself, because they determine whether the reported gain is scientifically defensible.
